\section{Defending the paper}\label{ch:defence}

Three internal reviews of the paper raised the objections below. Each is stated as a reviewer would put it, followed
by the honest answer, the evidence, and the experiment that would settle it with its status. An honest answer
concedes what the evidence does not show; a reviewer who finds a concession in the paper is less likely to find an
overclaim.

\subsection*{1. ``Isn't the whole gain just the readout?''}

\paragraph{Answer} Mostly, yes, and the paper says so. Without the readout for the others, FedECS-LT leads the
strongest other method by \R{main/cifar10/margin} and \R{main/cifar100/margin} points. With the readout for every
method of the GeFL-F pipeline it leads by \R{main/cifar10/er_margin} \R{main/cifar10/er_strongest_ci} and
\R{main/cifar100/er_margin} \R{main/cifar100/er_strongest_ci}, not significant after Holm's correction
($p=\R{main/cifar10/er_strongest_p}$ and $\R{main/cifar100/er_strongest_p}$), and by at most \R{nat/er_margin_hi} on the
natural datasets (\cref{sec:s-same-readout,sec:s-natural}). Balanced Softmax, which has no generator, comes within
\R{main/cifar10/er_diff/Balanced Softmax} points once it has the readout.

\paragraph{Why the readout is still part of the method} The readout follows from the diagnosis: \cref{sec:s-lastlayer}
shows that no per-client correction of the loss can give a last layer information about absent classes, and the
readout is the repair that does. The paper therefore defines FedECS-LT as ending with it, credits the family it
belongs to (Fed3R, FedNCM, FedCOF), and reports every method of the pipeline with it added, so the reader can see
exactly how much is the readout's. What the paper must not claim is that the generator side has a demonstrated effect
beyond the readout.

\paragraph{Settled by} The component ablation (RQ5) and the variants of GeFL-F, both pending (\cref{sec:s-pending}).

\subsection*{2. ``The name says exact class statistics, but half the components do not use them''}

\paragraph{Answer} Correct, and the paper now says which do. Three components use the secure statistics: the
anchored generator (or CCWA), the moment correction and the readout. Logit adjustment uses each client's own class
frequencies, consensus soft labels use the heads' predictions, and ten synthetic epochs are computation
(\cref{sec:s-recipe}). The name describes the first group and the readout, which carries most of the effect.

\subsection*{3. ``Is the weight-decay diagnosis causal?''}

\paragraph{Answer} Not shown. What is measured is that the rare classes' rows are several times shorter
(\cref{fig:s-rows}), as the contraction lemma predicts. Against a causal reading: GeFL-F's generated features are
recognised at \R{main/cifar10/GeFL-F/fidelity_all} on CIFAR-10-LT, \R{main/cifar10/GeFL-F/fidelity_head} on the head
classes and \R{main/cifar10/GeFL-F/fidelity_tail} on the tail classes, and at chance on CIFAR-100-LT
(\R{main/cifar100/GeFL-F/fidelity_all}). If the shrinkage were what loses the rare classes, the head classes, whose rows
are long, would be recognised; they are not. The row lengths also do not fall monotonically with class size. The
step-count heuristic $\kappa_c$ was not checked against the rows. The paper calls the shrinkage a measured property
of the generator, not a demonstrated cause.

\paragraph{Settled by} GeFL-F with decoupled weight decay alone (AdamW), pending. If it raises GeFL-F's tail recall by
more than the noise floor, the diagnosis gains a causal test; if not, the paper should say the shrinkage has no
measured consequence for accuracy.

\subsection*{4. ``The other diagnoses: do generated features without labels, and the moment correction, matter?''}

\paragraph{Answer} The label-free lemma is exact, but its size in the pipeline is not settled. The synthetic phase
does reduce the logit spread, and less for FedECS-LT than for GeFL-F; but FedDF, on GeFL-F's own generator, loses less
spread than FedECS-LT (\cref{sec:s-spread}), so the spread change is not decided by whether the features carry their
label. The moment correction brings the generated spreads to \R{main/mc_spread_lo} to
\R{main/mc_spread_hi} times the exact ones (\cref{sec:s-mc}), but its fidelity result is mixed: on CIFAR-10-LT it raises the mean fidelity from
\R{main/cifar10/FedECS-LT/fidelity_all} to \R{main/cifar10/FedECS-LT/fidelity_mc_all} and lowers the tail classes'
from \R{main/cifar10/FedECS-LT/fidelity_tail} to \R{main/cifar10/FedECS-LT/fidelity_mc_tail}; on CIFAR-100-LT the mean
is unchanged (\R{main/cifar100/FedECS-LT/fidelity_mc_all}). Its value for accuracy is unmeasured.

\paragraph{Settled by} The component ablation (``w/o moment correction'', ``w/o consensus soft labels''), pending.

\subsection*{5. ``Were the settings tuned on the test sets?''}

\paragraph{Answer} The development used the same splits (seeds 3--5) and the same test sets as the reported runs,
comparing about 36 configurations on CIFAR-10-LT and 17 on CIFAR-100-LT; the baselines were not tuned
(\cref{sec:s-development}). The CIFAR margins may therefore favour FedECS-LT. Two things limit the damage. The six
naturally imbalanced datasets were never used in development, and there the pattern repeats: large margins without the
readout for the others, small ones with it (\cref{sec:s-natural}). And every setting is the same on every dataset;
nothing was set per dataset.

\paragraph{Settled by} The run on new splits, seeds 6--8, pending.

\subsection*{6. ``Is the heads' advantage just more computation and ensemble distillation?''}

\paragraph{Answer} It may be. FedECS-LT trains ten synthetic epochs per round against one for GeFL-F and FedDF, with
consensus soft labels, which are a form of ensemble distillation; its training costs about two and a half times the
GFLOPs per round of GeFL-F (\cref{tab:s-cost}). The heads-only margin over FedDF, \R{main/cifar10/heads_margin} on
CIFAR-10-LT, is not significant ($p=\R{main/cifar10/heads_strongest_p}$) and has not been separated from this extra
computation.

\paragraph{Settled by} GeFL-F and FedDF with ten synthetic epochs, and FedDF with logit adjustment, pending.

\subsection*{7. ``With one client per architecture nothing is averaged; does the result hold with several?''}

\paragraph{Answer} Unknown. In every reported result each architecture has one client, so Stage III averages nothing,
and the per-client loss methods become purely local training. That is exactly the case in which the readout helps
most, because each last layer then fits one client's prior. With several clients per architecture, heads are
averaged and the readout may matter less.

\paragraph{Settled by} The 50-client runs (five per architecture), pending; their first attempt ran out of host
memory and was rebuilt.

\subsection*{8. ``If every client runs every head, why have different models at all?''}

\paragraph{Answer} The clients still \emph{train} different networks; but the consensus needs each client to run every
architecture's head on the generated pool every round, and the readout every body once. This is forward computation
only, and the paper states it as an assumption of the setting (\cref{ch:limits}). It does weaken the hardware
motivation of model heterogeneity: a client must be able to store and evaluate the largest head. It also means the
architectures are not secret.

\subsection*{9. ``Why is it worse than FedAvg with one shared model?''}

\paragraph{Answer} Because the setting removes the problem FedECS-LT addresses. With one shared model, FedAvg carries
every client's knowledge into every model, so a generator is not needed; and the last layer is trained by every
client, so the readout has nothing to repair and lowers accuracy (\R{t2/cifar10-r100/FedECS-LT/gain_er} points on
CIFAR-10-LT). FedECS-LT also trains its head on frozen features of unaugmented images while the baselines train end
to end with augmentation. \R{t2/cifar10-r100/above_list} are above it on CIFAR-10-LT (\cref{sec:s-shared}). The paper
reports this protocol as a test of scope, and the honest conclusion is that FedECS-LT is a method for clients with
different models only. Why the readout lowers accuracy here was not tested.

\subsection*{10. ``Is a small gain worth \R{cost/cifar10/comm_ratio_bsm} times the communication?''}

\paragraph{Answer} Against Balanced Softmax with the readout, on CIFAR-10-LT, the measured gain is
\R{main/cifar10/er_diff/Balanced Softmax} points, not resolved, for \R{cost/cifar10/comm_ratio_bsm} times the
communication and \R{cost/cifar10/time_ratio_bsm} times the time. On this evidence, a practitioner who only needs
balanced accuracy should use per-client Balanced Softmax plus the exact readout, which needs no generator. FedECS-LT's
extra cost is justified only if the pending ablations show that its generator side adds accuracy, or in settings,
such as several clients per architecture, where it may add more. Against GeFL-F, the base it extends, the overhead is
\R{cost/cifar10/comm_ratio} times the communication and \R{cost/cifar10/time_ratio} times the time.

\subsection*{11. ``What do the released statistics leak?''}

\paragraph{Answer} Secure aggregation hides each client's own sums, not the federation's. The server and every client
learn each class's mean and spread. For a class held by one client these are that client's statistics; for a small
class each image moves the mean noticeably; and with a two-convolution extractor a class mean is close to a mean image
at half resolution (\cref{sec:s-privacy}). The paper did not test an inversion of the means and has no differential
privacy, so it claims nothing beyond secure aggregation against an honest-but-curious server.

\paragraph{Settled by} An inversion test on the smallest classes, or a differentially private release with its
accuracy cost; neither is planned.

\subsection*{12. ``Why are the absolute accuracies so low?''}

\paragraph{Answer} Because the setting is GeFL-F's: a two-convolution extractor with 3 and 10 channels, frozen after
Stage I, and small heads, under a long tail with ten clients. GeFL's authors attribute GeFL-F's lower accuracy to
this extractor and these heads. In GeFL's own balanced setting the paper's earlier implementation of GeFL-F reaches
\R{repro/cifar10/GeFL-F/final}\% on CIFAR-10, against about \R{repro/cifar10/published}\% in GeFL's paper, so the
implementation is not the cause. A stronger or pre-trained extractor was not studied.

\paragraph{Settled by} The reference trained on all data pooled, planned for GeFL's 50-client setting and pending,
would give the ceiling of this extractor.

\subsection*{13. ``How is this different from FedCOF, Fed3R or CCVR?''}

\paragraph{Answer} The readout belongs to the same family as Fed3R (ridge regression), FedNCM (nearest class means) and
FedCOF (an LDA-type classifier); these also note that their sums do not depend on the partition. The difference is
where it is applied: per architecture, on bodies trained in the federation rather than pre-trained, with statistics
from every client, including clients of other architectures, which requires every client to run every body. CCVR
samples features from Gaussians with the pooled class moments to retrain a shared classifier; the moment correction
instead maps a learned generator's samples onto those moments. The pending ``w/ Gaussian sampler'' run compares the
two directly. The paper does not claim the readout or partition invariance as new.

\subsection*{14. ``Why only three seeds?''}

\paragraph{Answer} Compute: each configuration takes hours on a Kaggle GPU, and the study has several hundred runs.
Three seeds give little power. A sign-flip permutation test cannot go below $p=0.25$ with three pairs, so every
significant result rests on the $t$-test's normality assumption, and a difference of one to two points cannot be
resolved against a noise floor of \R{repeat/er/floor} (\cref{sec:s-stats-design}). The paper therefore reports
intervals, applies Holm's correction, calls small differences unresolved rather than absent, and adds three new seeds
for the main comparison (pending).

\subsection*{The position to defend}

The defensible claims are: three measured properties of GeFL-F under a long tail; a
method that, ending with an exact readout, has the highest balanced accuracy of the setting by a large margin; and the
finding that the readout, added to every method, carries most of that margin. Every other claim waits on a pending
run, and the paper says which.
